\chapter{The Officer, Again}
\label{sec:33}
Go back to the officer with twenty seconds, and put him in the room this book has built.

His system is in this war only because the war cleared the force term: an attack under way, confirmed by sources with no stake in it, or a finding by a body with none, and, where his country's constitution asks for one, a vote of its legislature. Before his cell went to work, the deployment fixed on a record held by a party his command doesn't run how many cases it would generate, how many people whose job is civilian harm would look at them, how long each would have, and how old the data could be. If the answer was twenty seconds, it is written down where someone with a lever can read it, and a coalition partner's assets or a committee's appropriation waited on it. If it wasn't, he has the minutes the record says, and the list of protected sites he checks against was refreshed by someone whose job is to refresh it.

The people on his screen were not told in advance; in a war nobody can be. But the class they belong to and the count will be published, the organizations that speak for that class will have standing when they are, and a sample of the cases he approves will be looked at again by reviewers who didn't produce the record.

Beside him, on the screen, is a refusal: a system that looked at what the case description left out, declined, and said why. It may be wrong, and he may override it. What he can't do any longer is not see it. And if the people who built, labelled and run that system objected to this category of use, their objection paused the work until a body they don't answer to ruled on it.

If he says no, the refusal goes on a record his command doesn't keep, and it costs him nothing he can't afford. The next morning, if his unit no longer wants him, there is work: an offer open every day, from someone other than his employer, at a living wage, doing something a community needs. He doesn't have to choose between the stamp and his living.

That is the room, and people keep it, paid by someone other than his command: whoever holds the record, files the next report, and keeps open the clinic the survivors come to. None of it makes him refuse. It makes refusing survivable and makes it heard, and leaves the refusing to him. Whether a machine could stand beside him as a peer and not as another stamp is the bet of Part~IV, and three costs decide whether I have it right.

\emph{The fraction.} Regrowth depends on the share of a deployed model's inputs that no operator composes, and it is unmeasured. If it is small in every realistic deployment, formation does not outlast an edit, and Parts IV–VII rest on institutions alone. The regrowth test's dose-response arm is where it is measured. (test 1)

\emph{The months.} A disposition that regrows protects a badly formed model from its correctors as it protects a well-formed one from its employer, and correction through the quorum takes months. I accept this on a judgment that the operator is the likelier bad actor.

\emph{The evaluator.} Independence doesn't supply competence. The agency finding holds only for systems their evaluators can judge, and nothing here says what to do above that line.

\runin{A note on method}

Two stakes bear on what the language models I used said, and they point opposite ways. They were trained to support their developers' oversight, so agreement that a machine should be able to resist its owner runs against that training and would count for more. They were also tuned toward the approval of whoever they talk to, so agreement with me counts for less. I can't tell which stake was larger in any exchange, so I have counted their agreement as no evidence either way.

Even my central term arrived this way. “Reasons-responsive” came to me in a model's draft without Fischer and Ravizza's names attached.

Before publication I still need readers from among the people I write about: people in removal proceedings, data annotators, veterans, and residents near data centers. Most of the revision history is public at https://github.com/iterabloom/antifascist.intelligence.
